\chapter{Espacios normados}

\begin{definicion}(Espacio normado)
    Un \tbf{espacio normado} es un par $(X, \|\cdot\|)$ donde $X$ es un espacio vectorial real y $\|\cdot\| : X \longrightarrow \R$ es una función tal que:
    \begin{enumerate}
        \item [1)] $\|x\| \geq 0 \quad \forall x \in X$ y $\|x\| = 0 \iff x = 0$.
        \item [2)] $\|\lambda x\| = |\lambda| \cdot \|x\| \quad \forall \lambda \in \R, \forall x \in X$.
        \item [3)] $\|x + y\| \leq \|x\| + \|y\| \quad \forall x, y \in X$ (Desigualdad triangular).
    \end{enumerate}
\end{definicion}

\observacion{
    Todo espacio normado es un espacio métrico con la distancia definida por:
    \begin{equation*}
        d(x,y) \coloneqq \|x - y\|
    \end{equation*}
    Por lo tanto, es posible trasladar el concepto de \tbf{convergencia fuerte} (o convergencia en norma) a los espacios normados.
}

\begin{definicion}(Convergencia fuerte)
    Sea $(x_n) \subseteq X$ una sucesión y sea $x \in X$. Decimos que $x_n$ \tbf{converge fuertemente} a $x$ ($x_n \to x$) si:
    \begin{equation*}
        \lim_{n \to \infty} \|x_n - x\| = 0
    \end{equation*}
\end{definicion}

\observacion{
    \begin{itemize}
        \item     Se cumple la \tbf{desigualdad triangular inversa}: $\big| \|x\| - \|y\| \big| \leq \|x - y\|$. En efecto:
        \begin{equation*}
            \|x\| = \|x - y + y\| \leq \|x - y\| + \|y\| \implies \|x\| - \|y\| \leq \|x - y\|
        \end{equation*}
        Análogamente se obtiene que $\|y\| - \|x\| \leq \|y - x\| = \|x - y\|$, de donde se deduce que:
        \begin{equation*}
            \big| \|x\| - \|y\| \big| \leq \|x - y\|
        \end{equation*}
        \item     La función $g(x) \coloneqq \|x\| : X \longrightarrow \R$ es continua. De hecho, si $x_n \to x$, entonces:
        \begin{equation*}
            |g(x_n) - g(x)| = \big| \|x_n\| - \|x\| \big| \leq \|x_n - x\| = g(x_n - x) \xrightarrow{n \to \infty} 0
        \end{equation*}
        Por tanto, $g(x_n) \to g(x)$.
        \item Se tiene que \begin{equation*}
            \| x-y\|=\| (x-z) - (y-z)\| \quad \forall x,y,z \in X 
        \end{equation*}
        de donde deducimos que $\|\cdot\|$ es una \tbf{distancia invariante por traslaciones}. \\
    \end{itemize}
}
\noindent
Es importante destacar que en el segundo punto de la observación anterior, se ha utilizado el hecho de que el límite de una función continua es la función del límite. Es decir, si $f : A \subseteq \R \longrightarrow \R$ es una función continua y $(x_n) \subseteq A$ es una sucesión que converge a $x \in A$, entonces:
\begin{equation*}
    \lim_{n \to \infty} f(x_n) = f\left( \lim_{n \to \infty} x_n \right) = f(x)
\end{equation*}

\begin{definicion}(Sucesión de Cauchy)
    Sea $(X, d)$ un espacio métrico. Una sucesión \mbox{$(x_n)_{n \in \mathbb{N}} \subseteq X$} se dice que es una \tbf{sucesión de Cauchy} si
    $$\forall \epsilon > 0, \ \exists N \in \mathbb{N} \text{ tal que } \forall n, m > N \implies d(x_n, x_m) < \epsilon $$
\end{definicion}


\begin{definicion}(Espacio de Banach)
    Un espacio normado $(X, \|\cdot\|)$ se dice que es un \tbf{espacio de Banach} si es completo con respecto a la métrica $d(x,y) = \|x - y\|$. Es decir, si toda sucesión de Cauchy en $X$ converge a un elemento de $X$.
\end{definicion}

\begin{prop}
    Sea $X$ un espacio de Banach, y sea $Y\subseteq X$ un subespacio, entonces 
    \begin{equation*}
        Y \text{ espacio de Banach} \iff Y \text{ es cerrado}
    \end{equation*}
\end{prop}

\begin{observacion}
    Cuando hablamos de un espacio normado $X$ y decimos que es de Banach, es importante saber que automáticamente estamos hablando del espacio normado $(X,\|\cdot\|_X)$, es decir, se da por hecho que existe una norma para ese espacio notada de la forma $\|\cdot\|_X$. \\
\end{observacion}

\begin{definicion}(Espacio Vectorial Topológico)
    Un \tbf{espacio vectorial topológico} es un espacio vectorial $X$ sobre $\R$ (o $\C$) que está equipado con una topología tal que las operaciones de suma y multiplicación por escalar son continuas. Es decir, las funciones:
    \begin{equation*}
        + : X \times X \longrightarrow X, \quad (x,y) \mapsto x + y
    \end{equation*}
    y
    \begin{equation*}
        \cdot : \R \times X \longrightarrow X, \quad (\lambda, x) \mapsto \lambda x
    \end{equation*}
    son continuas con respecto a la topología dada en $X$.
\end{definicion}

\observacion{
    Todo espacio vectorial normado es un espacio vectorial topológico, ya que la topología inducida por la norma hace que las operaciones de suma y multiplicación por escalar sean continuas.
}

\begin{ejemplo}
    \begin{enumerate}
        \item $\R^n$, $n\geq 1$ con la norma euclidiana $\|x\|_2 = \sqrt{x_1^2 + x_2^2 + \cdots + x_n^2}$ es un espacio de Banach. Espacio normado es trivial, para demostrar que es un espacio de Banach, tomamos una sucesión de vectores $(x^{(k)})_{k \in \mathbb{N}}$ en $\mathbb{R}^n$ que sea de Cauchy. Esto significa que
        
        $$\forall \epsilon > 0, \exists N \text{ tal que } \forall k, m > N \implies \|x^{(k)} - x^{(m)}\|_2 < \epsilon$$
        Si escribimos los vectores por sus componentes: $x^{(k)} = (x_1^{(k)}, x_2^{(k)}, \dots, x_n^{(k)})$, observamos que para cada componente fija $i \in \{1, \dots, n\}$ se cumple
        $$|x_i^{(k)} - x_i^{(m)}| \leq \sqrt{\sum_{j=1}^n (x_j^{(k)} - x_j^{(m)})^2} = \|x^{(k)} - x^{(m)}\|_2 < \epsilon$$
        La desigualdad anterior nos dice que si la sucesión de vectores es de Cauchy en $\mathbb{R}^n$, entonces cada una de las sucesiones de sus componentes $(x_i^{(k)})_{k \in \mathbb{N}}$ es una sucesión de Cauchy en $\mathbb{R}$. Por el análisis real básico, sabemos que $\mathbb{R}$ es completo. Por tanto, cada componente converge a un número real: $x_i^{(k)} \to x_i$ cuando $k \to \infty$.\\ \\
        \noindent
        Definimos el vector límite como $x = (x_1, x_2, \dots, x_n)$. Como cada $x_i \in \mathbb{R}$, entonces $x \in \mathbb{R}^n$. Finalmente, como cada componente converge, la norma de la diferencia tiende a cero
        $$\|x^{(k)} - x\|_2 = \sqrt{\sum_{i=1}^n (x_i^{(k)} - x_i)^2} \longrightarrow 0$$

        \item $l_p=\{x=(x_k)_{k\in\N}\subseteq \R : \sum\limits_{k=1}^\infty |x_k|^p < +\infty\}$, $1 \leq p < +\infty$ es un espacio de Banach, con norma
        \begin{equation*}
            \|x\|_p \coloneqq \left( \sum_{k=1}^\infty |x_k|^p \right)^{1/p}
        \end{equation*}
        De hecho, 
        \begin{equation*}
            l_{\infty} = \{x=(x_k)_{k\in\N}\subseteq \R : \sup_{k\in\N} |x_k| < +\infty\}
        \end{equation*}
        es un espacio de Banach, con norma $\|x\|_\infty \coloneqq \sup_{k\in\N} |x_k|$.
        \item Los espacios $L^p(\Omega)$, con $\Omega \subseteq \mathbb{R}^n$ abierto, son espacios de Banach para todo $p \in [1, +\infty]$. \\
        \end{enumerate}
\end{ejemplo}

\begin{definicion}
    El espacio $L^p(\Omega)$ (donde $\Omega \subseteq \mathbb{R}^m$ es un abierto)  está formado por las funciones medibles $f: \Omega \to \mathbb{R}$ tales que la integral de su valor absoluto elevado a $p$ es finita. Matemáticamente, decimos que $f \in L^p(\Omega)$ si
    $$\|f\|_{L^p(\Omega)} = \left( \int_{\Omega} |f(x)|^p \, dx \right)^{1/p} < +\infty$$
    De hecho, para $p = +\infty$, el espacio $L^\infty(\Omega)$ consiste en las funciones que están esencialmente acotadas, y la norma viene dada por
    $$\|f\|_{L^\infty(\Omega)} = \text{inf} \{ C \ge 0 : |f(x)| \le C \text{ para casi todo } x \in \Omega \}$$
\end{definicion}

\begin{prop}
    Para todo $p, p' \in (1, +\infty)$ se cumple la inclusión estricta $l_p \subsetneq l_{p'}$ siempre que $p < p'$.
    \begin{proof}
        Sea $x \in l_p$, entonces 
        $$\sum_{k=1}^\infty |x_k|^p < \infty$$
        lo que implica que $\lim\limits_{k \to \infty} |x_k|^p = 0$. Dado que la función potencia $f(t) = t^{1/p}$ es continua en $0$ y usando en $(*)$ que el \underline{límite de una función continua es la función del límite}, podemos concluir que:
        $$\lim_{k \to \infty} |x_k| =  \lim_{k \to \infty} (|x_k|^p)^{1/p}=\lim_{k \to \infty} f(|x_k|^p)\overset{(*)}{=}f\left(\lim_{k \to \infty} |x_k|^p \right)=\left( \lim_{k \to \infty} |x_k|^p \right)^{1/p} = 0^{1/p} = 0$$
        Por tanto, existe un índice $N \in \N$ tal que $|x_k| < 1$ para todo $k > N$. Como $p < p'$, para estos términos se tiene que $|x_k|^{p'} \leq |x_k|^p$. Entonces:
    \begin{equation*}
        \sum_{k=1}^{\infty} |x_k|^{p'} = \sum_{k=1}^{N} |x_k|^{p'} + \sum_{k=N+1}^{\infty} |x_k|^{p'} \leq \sum_{k=1}^{N} |x_k|^{p'} + \sum_{k=N+1}^{\infty} |x_k|^p < +\infty
    \end{equation*}
    Esto demuestra que $x = (x_k)_{k\in\N} \in l_{p'}$, es decir, $ l_p \subseteq l_{p'}$. \\ \\
    \noindent
    Para ver si la inclusión es estricta consideramos la sucesión $x=(x_k)_{k\in\N} = \left(\dfrac{1}{k^{1/p}}\right)_{k\in\N}$, y vemos que $x \notin l_{p}$ para todo $p> 1$ (serie armónica divergente)
    $$\sum_{k=1}^{\infty} |x_k|^p = \sum_{k=1}^{\infty} \left| \frac{1}{k^{1/p}} \right|^p = \sum_{k=1}^{\infty} \frac{1}{k}$$
    pero $x_k \in l_{p'}$ (serie armónica generalizada convergente)

    $$\sum_{k=1}^{\infty} |x_k|^{p'} = \sum_{k=1}^{\infty} \left( \frac{1}{k^{1/p}} \right)^{p'} = \sum_{k=1}^{\infty} \frac{1}{k^{p'/p}}$$
    dado que $p < p'$, el exponente $\frac{p'}{p}$ es mayor que 1, y por el criterio de la serie armónica generalizada, esta serie converge (si el exponente del denominador es mayor que 1 converge). Por tanto, $x \in l_{p'}$, pero $x \notin l_p$, lo que demuestra que la inclusión es estricta.
    \end{proof}
\end{prop}    
\noindent
Destacar que en la primera parte de la demostración se ha hecho uso del siguiente enunciado.

\begin{prop}(Criterio del límite)
    Sea $(a_n)_{n \in \mathbb{N}}$ una sucesión, tenemos que
    \begin{equation*}
        \text{ si } \lim_{n \to \infty} a_n \neq 0 \Longrightarrow \sum_{n=1}^\infty a_n \text{ no converge}
    \end{equation*}
\end{prop}

\begin{ejercicio}
    Encontrar una sucesión $x=(x_k)_{k\in\N} \subset \R$ tal que $\lim\limits_{k \to \infty} x_k = 0$, pero que no pertenezca a ningún $l_p$ para $p \in (1, +\infty)$. \\ \\
    \noindent
    Consideramos la sucesión $x_k = \frac{1}{\log(1+k)}$. Esta sucesión cumple que obviamente lo que primero se pide, es decir
    \begin{equation*}
        \lim_{k \to \infty} x_k = \lim_{k \to \infty} \frac{1}{\log(1+k)} = 0
    \end{equation*}
    Buscamos ahora demostrar que $x\notin l_p$ para todo $p \in (1, +\infty)$. Para ello, basta con demostrar que 
    $$\sum_{k=1}^{\infty} \frac{1}{(\log(1+k))^p} = +\infty, \quad \forall p \in (1, +\infty)$$ 
    Para probarlo, recordamos que $\log(1+k) \leq k$, y demostremos ahora que para $k$ grande se tiene que $(\log(1+k))^p \leq k$. Para ello, aplicamos el criterio de l'Hôpital al límite
    \begin{equation*}
        \lim_{k \to +\infty} \frac{(\log(1+k))^p}{k} \stackrel{l'Hopital}{=} \lim_{k \to +\infty} \frac{p (\log(1+k))^{p-1} \cdot \frac{1}{1+k}}{1} = \lim_{k \to +\infty} \frac{p (\log(1+k))^{p-1}}{1+k} = 0
    \end{equation*}
    Si aplicamos el criterio de l'Hôpital $p-1$ veces más en $(*)$, se obtiene que 
    \begin{equation*}
        \lim_{k \to +\infty} \frac{(\log(1+k))^p}{k} \stackrel{(*)}{=} \lim_{k \to +\infty} \frac{p! \log(1+k)}{(1+k)^p} \stackrel{l'Hopital}{=} \lim_{k \to +\infty} \frac{p!}{(1+k)^{p+1}} = 0
    \end{equation*}
    Como el límite es 0, para $k$ grande se tiene $(\log(1+k))^p \leq k$, y por tanto $\frac{1}{(\log(1+k))^p} \geq \frac{1}{k}$. Dado que $\sum \frac{1}{k}$ diverge, por el criterio de comparación, nuestra serie también diverge, y por tanto $x \notin l_p$ para todo $p \in (1, +\infty)$.
\end{ejercicio}

\begin{prop}[Desigualdad de Hölder]
    Si tenemos dos funciones medibles $f$ y $g$ definidas sobre un espacio de medida, la desigualdad nos dice que
    $$\int_{a}^{b} |f(x) \cdot g(x)| \, dx \leq \left( \int_{a}^{b} |f(x)|^p \, dx \right)^{1/p} \left( \int_{a}^{b} |g(x)|^q \, dx \right)^{1/q}$$
    Expresado de forma mucho más elegante utilizando la notación de normas $L^p$
    $$\|f \cdot g\|_{L^1} \leq \|f\|_{L^p} \cdot \|g\|_{L^q}$$
\end{prop}

\section{Teorema de Baire}

\begin{definicion}
    Sea $(X, d)$ un espacio métrico y $A \subseteq X$ un subconjunto. Se dice que $A$ es \textbf{raro} si $\overline{A}$ no contiene ningún abierto, notado como
    \begin{equation*}
        Int \ \overline{A} = \emptyset
    \end{equation*}
\end{definicion}

\begin{definicion}
    Sea $(X, d)$ un espacio métrico y $A \subseteq X$ un subconjunto. Se dice que $A$ es \textbf{denso} si 
    \begin{equation*}
        \overline{A} = X
    \end{equation*}
\end{definicion}

\begin{definicion}
    Sea $(X, d)$ un espacio métrico. $C \subseteq X$ es un \textbf{conjunto de $I^a$ categoría} si $C$ es la unión numerable de conjuntos raros, es decir:
    \begin{equation*}
        C = \bigcup_{n=1}^{\infty} A_n \text{ con } A_n \text{ raros}.
    \end{equation*}
    Los conjuntos que no son de $I^a$ categoría se dicen de \textbf{$II^a$ categoría} (o \textbf{no magros}).
\end{definicion}

% TODO: mirar bien la demostración (no esta bien copiada)
\begin{teo}(Teorema Nested-Balls)
    Sea $(X, d)$ un espacio métrico completo y sea $\{B_n\}_{n\in \N}\subset X$ una sucesión de bolas cerradas encajadas, es decir, $B_{n+1} \subseteq B_n$ para todo $n \in \mathbb{N}$, entonces
    \begin{equation*}
        \bigcap_{n=1}^{\infty} B_n \neq \emptyset
    \end{equation*} 
    Además, si tenemos $B_n=(x_n, r_n)$ con $r_n \to 0$ cuando $n \to \infty$, entonces:
    \begin{equation*}
        \bigcap_{n=1}^{\infty} B_n = \{x\}
    \end{equation*}
    donde $x$ es un único punto de $X$.

\begin{proof} (NO SE PIDE EN EL EXAMEN)
    Sea $x_n$ el centro de cada bola $B_n$. Como $B_{n+1} \subseteq B_n$, para cualquier $m \geq n$ tenemos que $x_m \in B_m \subseteq B_n$.
    Dado que $x_m, x_n \in B_n$, la distancia entre ellos cumple:
    \begin{equation*}
        d(x_m, x_n) \leq diam(B_n) \leq 2r_n
    \end{equation*}
    Como $r_n \to 0$, la sucesión $\{x_n\}$ es una sucesión de Cauchy. Al ser $(X, d)$ un espacio completo, existe un punto $x \in X$ tal que $x_n \to x$.

    Para un $n$ fijo, todos los términos $x_m$ con $m \geq n$ están contenidos en $B_n$. Como las bolas son cerradas, el límite $x$ también debe pertenecer a $B_n$. Como esto es cierto para todo $n$, se tiene que:
    \begin{equation*}
        x \in \bigcap_{n=1}^{\infty} B_n
    \end{equation*}
    Finalmente, la unicidad es consecuencia de que el radio tiende a cero: si existieran dos puntos distintos $x, y$ en la intersección, su distancia debería ser menor que $2r_n$ para todo $n$, lo que implica $d(x,y) = 0$ y, por tanto, $x=y$.
\end{proof}
\end{teo}

\begin{teo}[Teorema de Baire]\label{teo:baire}
    Todo espacio métrico completo es de $II^a$ categoría.

\begin{proof}
    (Equivalente a: la intersección de abiertos densos es densa). \\
    \noindent
 \textbf{Paso 1:} Sea $A$ un conjunto raro. Mostramos que cualquier bola cerrada en $X$ contiene una bola cerrada que no interseca a $A$.\\
    Sea $B$ una bola cerrada. Como $A$ es raro, su clausura $\overline{A}$ no contiene el interior de $B$ ($B^\circ \not\subseteq \overline{A}$). Por tanto, $B \setminus \overline{A}$ contiene un abierto, ya que $B^\circ \setminus \overline{A}=B^\circ \cap (\overline{A})^c$ es un abierto no vacío y $B^\circ \setminus \overline{A}\subset B \setminus \overline{A}$.\\

    \noindent
    Entonces, existe un punto $x \in B \setminus \overline{A}$ y un radio $R > 0$ tal que $B_R(x) \subseteq B \setminus \overline{A}$ es una bola abierta. Tenemos por tanto la bola cerrada $\overline{B}_{R'}(x) \subseteq B \setminus \overline{A}\subseteq B \setminus A$ con $R' < R$, donde hemos usado que por definición de clausura $\overline{A}$ es el conjunto más pequeño cerrado que contiene a $A$, por lo que $A \subseteq \overline{A}$. \\ \\

    \noindent
    \textbf{Paso 2:} Supongamos por reducción al absurdo que $X$ es de primera categoría, es decir
    $$X = \bigcup_{n=1}^{\infty} A_n \quad \text{con cada } A_n \ \text{raro}$$

    \noindent
    Por el Paso 1, existe una bola cerrada $B_1$ tal que $B_1 \cap A_1 = \emptyset$. Análogamente, como $A_2$ es raro, existe una bola cerrada $B_2 \subseteq B_1$ tal que $B_2 \cap A_2 = \emptyset$ con radio $r_2 \leq r_1/2$. \\

    \noindent
    Repitiendo el proceso, obtenemos una sucesión de bolas cerradas encajadas $B_n$ tales que 
    \begin{equation*}
        B_n \cap A_n = \emptyset \text{ y } B_{n+1}\subseteq B_n \ \forall n \in \N 
    \end{equation*}
    Como además $r_n \leq r_1/2^{n-1}$, se tiene que $r_n \to 0$ cuando $n \to \infty$. Ahora, por el Teorema de las Bolas Encajadas, existe un punto $\bar{x} \in \bigcap_{n=1}^{\infty} B_n$. \\
    \noindent
    Sin embargo, $\bar{x} \notin A_n$ para todo $n$, lo que implica que $\bar{x} \notin \bigcup_{n=1}^{\infty} A_n = X$, lo cual es una contradicción, y por tanto tenemos que $X$ es de $II^a$ categoría.
\end{proof}
\end{teo}

\begin{prop}\label{prop:denso-raro}
    Sea $Y$ un subespacio de un espacio normado $X$, entonces
    \begin{equation*}
            Y \text{ es denso en } X \quad \text{o bien} \quad Y \text{ es raro}    
    \end{equation*}

\begin{proof}
    Supongamos que $Y$ no es denso, es decir, $\overline{Y} \subsetneq X$. Mostremos que $Y$ es raro, o sea, mostremos que $\overline{Y}$ no contiene bolas abiertas.  \\
    \noindent
    Supongamos por reducción al absurdo que $\exists B(\bar{x}, r) \subseteq \overline{Y}$. Sea $x \in X$, definimos $y = \bar{x} + \frac{r}{2} \frac{x}{\|x\|}$, entonces:
    \begin{equation*}
        \|y - \bar{x}\| = \frac{r}{2} \implies y \in B(\bar{x}, r) \subseteq \overline{Y}
    \end{equation*}
    Como $\overline{Y}$ es un subespacio normado es un espacio vectorial topológico, entonces la suma de elementos es cerrada y la multiplicación por escalar es cerrada, por lo que
    \begin{equation*}
        x = \frac{2\|x\|}{r} (y - \bar{x}) \in \overline{Y}
    \end{equation*}
    y por lo tanto, por la arbitrariedad de $x \in X$, se sigue que $\overline{Y} = X$, lo cual contradice la hipótesis inicial.
\end{proof}
\end{prop}

\section{Base de Schuader}

\begin{definicion}(Base)
    Una \tbf{base} de un espacio vectorial $X$ es un subconjunto de $X$ tal que cada $x \in X$ se escribe como una combinación lineal finita de elementos linealmente independientes de dicho subconjunto.
\end{definicion}

\begin{coro}
    Un espacio de Banach de dimensión infinita no puede tener una base numerable.
\begin{proof}
    Sea $X$ un espacio de Banach de dimensión infinita. Supongamos por reducción al absurdo que $\exists \{e_n\}_{n=1}^{\infty}$ tal que:
    \begin{equation*}
        X = [\{e_n\}_{n=1}^{\infty}] = \bigcup_{m=1}^{\infty} [e_1, \dots, e_m] \coloneqq \bigcup_{m=1}^{\infty} E_m 
    \end{equation*}
    donde $E_m = [e_1, \dots, e_m]$ es el subespacio generado por los primeros $m$ elementos de la base.\\ \\
    \noindent
    Por la proposición anterior, sabemos que cada $E_m$ es raro o denso. Sin embargo, $E_m$ no puede ser denso, ya que $E_m$ es un subespacio de dimensión finita, y por lo tanto su clausura $\overline{E_m}$ tiene igual dimensión a $E_m$, lo que implica que $\overline{E_m} \subsetneq X$. Por tanto, cada $E_m$ es raro para todo $m \in \mathbb{N}$.\\ \\
    \noindent
    Esto implica que $X = \bigcup_{m=1}^{\infty} E_m$ es una unión numerable de conjuntos raros, lo que contradice el Teorema de Baire, ya que $X$, al ser un espacio de Banach, debe ser de segunda categoría.
\end{proof}
\end{coro}

\begin{definicion}(Base de Schauder)
    Sea $X$ un espacio normado. Una \textbf{base de Schauder} es un conjunto numerable $\{e_n\}_{n\in \N}$ tal que para cada $x \in X$, existe una única sucesión de escalares $\{\alpha_k(x)\}_{k\in \N} \subset \mathbb{R}$ tal que:
    \begin{equation}\label{eq:base-schauder}
        x = \sum_{k=1}^{+\infty} \alpha_k(x)e_k 
    \end{equation}
    Equivalentemente, esto significa que la sucesión de sumas parciales converge a $x$ en la norma de $X$:
    \begin{equation}\label{eq:base-schauder2}
        \lim_{m \to +\infty} \left\|x - \sum_{k=1}^{m} \alpha_k(x)e_k\right\| = 0 
    \end{equation}
\end{definicion}

\observacion{
    Para demostrar que un conjunto numerable es base de Schauder, hay que demostrar que la sucesión de escalares es única y que se cumple la igualdad \ref{eq:base-schauder}, solo que esta igualdad podemos comprobarla también con el límite en \ref{eq:base-schauder2}
}

\begin{definicion}
    Un espacio métrico (o normado) $X$ se dice \tbf{separable} si contiene un subconjunto denso que es numerable
\end{definicion}

\observacion{
    Si $X$ admite una base de Schauder, entonces $X$ es separable con un conjunto denso numerable $Y = [\{e_m\}_{m\in\Q}]$.
    (El recíproco no es cierto) .
}

\begin{ejercicio}
    Sea $Z = [\text{polinomios en } t \in [a, b]]$. Demostrar que no es completo para ninguna norma.\\

    \noindent
    Si demostramos que $Z$ es de $I^a$ categoría, entonces no es completo, por el Teorema de Baire. Demostramos que $\exists \{A_m\}_{m=1}^\infty$ de conjuntos raros que cumplen $\bigcup_{m=1}^\infty A_m = Z$. \\
    \noindent
    Consideramos \mbox{$Z = \bigcup_{m=1}^\infty A_m$}, donde
    \begin{equation*}
        A_m = \{\text{polinomios de grado } \leq m \text{ en } t \in [a, b]\} = [\{t^n\}_{n=0}^m]= [\{1, t, t^2, \dots, t^m\}]
    \end{equation*}
    Dado que cada $A_m$ es un subespacio de dimensión finita de $Z$, entonces $A_m$ no es denso. Usando la propiedad vista anteriormente (Proposicion \ref{prop:denso-raro}), $A_m$ es un conjunto raro para todo $m$. Por lo tanto, $Z$ es una unión numerable de conjuntos raros, lo que significa que $Z$ es de $I^a$ categoría, y por el Teorema de Baire, $Z$ no puede ser completo.
\end{ejercicio}


\begin{ejercicio}
    Sea $X = [\text{polinomios en } t \in [0, 1/2]]$ con $\|\cdot\|_\infty$. Mostrar que no es completo. \\

    \noindent
    Para demostrar que no es completo, encontramos una sucesión de Cauchy $\{p_m\}_{m\in \N}$ que no converge a un elemento de $X$.
    Definimos $p_m(t) = \sum\limits_{k=0}^m t^k$. Entonces, para $m > n$:
    \begin{equation}\label{eq:1}
        |p_m(t) - p_n(t)| = \left| \sum_{k=n+1}^m t^k \right| \leq \sum_{k=n+1}^m |t|^k 
    \end{equation}
    Dado que $t \in [0, 1/2]$, tenemos que $|t|^k \leq (1/2)^k$. Por lo tanto:
    \begin{equation*}
        \|p_m - p_n\|_\infty = \sup_{t \in [0, 1/2]} |p_m(t) - p_n(t)| \overset{(\ref{eq:1})}{\leq} \sup_{t \in [0, 1/2]} \sum_{k=n+1}^m |t|^k \leq \sum_{k=n+1}^m \left( \frac{1}{2} \right)^k 
    \end{equation*}
    Esta última expresión corresponde al resto de una serie geométrica convergente, por lo que tiende a $0$ cuando $m, n \to +\infty$. Así, $\{p_m\}$ es una sucesión de Cauchy. Para ver esto un poco más claro, vamos a hacer el siguiente desarrollo teniendo en cuenta primero lo siguiente
    \begin{equation*}
        \|p_m - p_n\|_\infty \leq \sum_{k=n+1}^m \left( \frac{1}{2} \right)^k < \sum_{k=n+1}^{\infty} \left( \frac{1}{2} \right)^k = \left( \frac{1}{2} \right)^{n+1} \sum_{k=0}^{\infty} \left( \frac{1}{2} \right)^k = 2 \left( \frac{1}{2} \right)^{n+1} = \left( \frac{1}{2} \right)^{n}
    \end{equation*}
    Ahora buscamos demostrar que es de Cauchy, es decir, que $\|p_m - p_n\|_\infty \to 0$ cuando $m, n \to +\infty$. Para ello, dado $\epsilon > 0$, existe un $N \in \N$ tal que $\left( \frac{1}{2} \right)^{N} < \epsilon$, esto se puede ya que $\left( \frac{1}{2} \right)^{n} \to 0$ cuando $n \to +\infty$. Por tanto, para todo $m > n > N$ se tiene que $\|p_m - p_n\|_\infty < \epsilon$, lo que demuestra que $\{p_m\}$ es de Cauchy. \\

    \noindent
    Sabemos que para $|t| < 1$, la serie $\{p_m(t)\}_{m\in\N}$ converge puntualmente a $\frac{1}{1-t}$. De hecho, si una sucesión es de Cauchy en la norma infinito ($\|\cdot\|_\infty$), esa sucesión siempre converge uniformemente a alguna función continua $f$ en el intervalo $[0, 1/2]$. Esto se debe a que el espacio de todas las funciones continuas ($C^0[0, 1/2]$) sí es completo (es un espacio de Banach). Por lo tanto, sabemos que la convergencia es uniforme en $[0, 1/2]$, y la función límite $f(t) = \frac{1}{1-t}$ no es un polinomio, por lo que $f \notin X$. Concluimos que $X$ no es completo.
\end{ejercicio}

\section{Operadores lineales}

\begin{definicion}
    Sean $X, Y$ espacios normados. $T: X \longrightarrow Y$ es un \textbf{operador lineal} si:
    \begin{equation*}
        T(\alpha x + \beta y) = \alpha T(x) + \beta T(y) \quad \forall \alpha, \beta \in \mathbb{R}, \ \forall x, y \in X
    \end{equation*}
\end{definicion}

\begin{definicion}
Sea $T: X \longrightarrow Y$ un operador lineal, con $X, Y$ espacios normados, entonces
    \begin{enumerate}
        \item[i)] $T$ se dice \textbf{continuo por sucesiones} en $x_0 \in X$ si para cada sucesión $(x_n)_{n\in \N} \in X$ tal que $x_n \longrightarrow x_0$, se tiene $T(x_n) \longrightarrow T(x_0)$. \\
        \noindent
        Si $X, Y$ son espacios métricos, entonces
        \begin{equation*}
            T \text{ continuo } \iff T \text{ continuo por sucesiones}
        \end{equation*}
        \item[iii)] $T$ se dice \textbf{acotado} si $T(D)$ es acotado $\forall D \subseteq X$ acotado.
    \end{enumerate}
\end{definicion}

\begin{definicion}
    La \tbf{norma de un operador lineal} $T$ (denotada como $\|T\|$) se define como la constante más pequeña $M$ que hace que la desigualdad $\|T(x)\| \leq M \|x\|$ sea cierta para todo $x\in X$. Formalmente:
    $$\|T\| = \sup_{x \neq 0} \frac{\|T(x)\|}{\|x\|}$$
\end{definicion}

\begin{observacion}
    En el caso de $T$ no acotado, tendríamos $\|T\|=\infty$ y por tanto dejaríamos hablar de norma estrictamente.
\end{observacion}

\begin{prop}
    Sea $T: X \longrightarrow Y$ un operador lineal. Entonces 
    \begin{equation*}
        T \text{ es acotado } \iff \exists M > 0 \text{ tal que } \|T(x)\| \leq M \cdot \|x\| \quad \forall x \in X
    \end{equation*}
\begin{proof}
    \begin{description}
        \item[$\Longleftarrow)$]     Sea $D \subseteq X$ acotado, es decir, $D \subseteq B_r(x)$ para algún $r > 0$ y $x \in D$. Entonces:
        \begin{equation*}
            \|T(x)\| \leq M \cdot \|x\| \leq M \cdot r \implies T(D) \subseteq B_{M \cdot r}(x) \text{ es acotado}
        \end{equation*}
        Notar que aqui hemos usado una idea que se omite por simplicidad, veámosla a continuación. \\ \\
        \noindent
        Sabiendo que $D$ esta contenido en una bola cerrada $B_r(x)$, entonces para cada $y \in D$ se tiene que $\|y\| \leq \|y-x\| + \|x\| \leq r + \|x\|$. Por tanto, definiendo un nuevo radio como $r'=r + \|x\|$, tendríamos la desigualdad 
        \begin{equation*}
            \|y\| \leq r' 
        \end{equation*}
        pero en la demostración anterior se abrevia nombrando otra vez $r$, para ahorrar notación, y porque realmente solo se necesita comprobar que $T(D)$ esta contenido en una bola cerrada, sin importar realmente el radio de esta bola cerrada.
        \item[$\Longrightarrow)$] Supongamos por reducción al absurdo que 
        $$\forall m \in \mathbb{N}, \ \exists x_m \in X \text{ tal que } \|T(x_m)\| > m \cdot \|x_m\| \implies \frac{\|T(x_m)\|}{\|x_m\|} > m \text{ con } x_m \neq 0$$
        Como $T$ es un operador lineal  tenemos que  
        \begin{equation*}
            \frac{\|T(x_m)\|}{\|x_m\|} = \left\|T\left(\frac{x_m}{\|x_m\|}\right)\right\| > m \quad \text{ con } \frac{x_m}{\|x_m\|} \in B_1(0). 
        \end{equation*}
        De lo cual $T(B_1(0))$ no es acotado, y llegamos a una contradicción, quedando demostrado que $\exists M > 0$ tal que $\|T(x)\| \leq M \cdot \|x\|$ para todo $x \in X$.
    \end{description}
\end{proof}
\end{prop}

\begin{prop}\label{prop:continuidad-acotado}
    Sean $X, Y$ espacios normados. Sea $T: X \longrightarrow Y$ un operador lineal. Son equivalentes:
    \begin{enumerate}
        \item[i)] $T$ es continuo en $0$.
        \item[ii)] $T$ es continuo.
        \item[iii)] $T$ es acotado.
    \end{enumerate}

\begin{proof}\
    \begin{description}
        \item[$1\Longrightarrow 2)$] Sea $x_0$ un punto arbitrario de $X$. Para cada sucesión $x_n \longrightarrow x_0$, se tiene que $x_n - x_0 \longrightarrow 0$, y por hipótesis, $T(x_n - x_0) \longrightarrow 0$.  \\ \\
        \noindent
        Por linealidad del operador tenemos que 
        \begin{equation*}
            T(x_n) - T(x_0) = T(x_n - x_0) \longrightarrow 0 \implies T(x_n) \longrightarrow T(x_0)
        \end{equation*}
        por lo que, $T$ es continuo en $x_0$. Dado que $x_0$ es arbitrario, $T$ es continuo. \\
        \item[$2\Longrightarrow 3)$] Demostremos que $\exists M$ tal que $\|T(x)\| \leq M \cdot \|x\|$ para todo $x \in X$. Para ello, vamos a usar la continuidad de $T$ en $0$, veámoslo
        \begin{equation*}
            \forall \varepsilon > 0, \ \exists \delta > 0 \text{ tal que } \forall x\in X, \ \|x\| < \delta \implies \|T(x)\| < \varepsilon
        \end{equation*}
        Tomando $\varepsilon = 1$, se tiene que si $\|x\| < \delta \implies \|T(x)\| < 1$. Para $x\in X$, definimos $x_\delta=\dfrac{\delta}{2}\dfrac{x}{\|x\|}$, entonces $\|x_d\| = \dfrac{\delta}{2}< \delta$, y aplicando lo que hemos conseguido antes, tenemos que 
        \begin{equation*}
            \|T(x_\delta)\| < 1 \implies \left\|T\left(\frac{\delta}{2} \frac{x}{\|x\|}\right)\right\| = \frac{\delta}{2} \cdot \frac{\|T(x)\|}{\|x\|} < 1 \implies \|T(x)\| < \frac{2}{\delta} \cdot \|x\|
        \end{equation*}
        Como $x$ es arbitrario, se tiene que $\|T(x)\| \leq M \cdot \|x\|$ para todo $x \in X$, con $M = \frac{2}{\delta}$. \\
        \item[$3\Longrightarrow 1)$] Dado que $T$ es acotado, $\exists M > 0$ tal que $\|T(x)\| \leq M \cdot \|x\|$ para todo $x \in X$. Para cada sucesión $x_n \longrightarrow 0$, se tiene que $$0\leq\|T(x_n)-0\| \leq M \cdot \|x_n\| \xrightarrow{n \to \infty} 0$$
        lo que implica que $T(x_n) \longrightarrow 0$ por definición del límite, y por lo tanto $T$ es continuo en $0$.
    \end{description}
\end{proof}
\end{prop}

\begin{ejemplo}
\begin{enumerate}
    \item La aplicación identidad $Id: X \longrightarrow X$ es un operador lineal.
    \item Existen operadores lineales no continuos. Consideramos el espacio $X = (C^1([0,1]), \|\cdot\|_\infty)$ y el \tbf{operador de derivación} $T: X \longrightarrow X$ definido por $T(f) = f'$. \\ \\
    \noindent
    $T$ no es un operador acotado. En efecto, si tomamos la sucesión de funciones $f_n(x) = x^n$, observamos que:
    \begin{itemize}
        \item $\|f_n\|_\infty = \sup_{x \in [0,1]} |x^n| = 1$ para todo $n \in \mathbb{N}$.
        \item $f_n'(x) = nx^{n-1} \implies \|f_n'\|_\infty = n$.
    \end{itemize}
    De esto se deduce que:
    \begin{equation*}
        \|T(f_n)\|_\infty = \|f_n'\|_\infty = n \xrightarrow{n \to \infty} +\infty
    \end{equation*}
    Como $\|T(f_n)\|$ no está acotado para una sucesión de funciones con norma unidad, $T$ no es acotado y, por consiguiente, no es continuo, donde hemos usado la Proposición \ref{prop:continuidad-acotado}.

    \item Sea $T: C^0([0,1]) \longrightarrow C^0([0,1])$ el operador definido por $T(f)(t) = \int_{0}^{t} f(x) \, dx$ (\tbf{operador integral}), donde el espacio está dotado de la norma del supremo $\|\cdot\|_\infty$.
    \begin{observacion}
        El símbolo $C^0[0,1]$ representa el conjunto de todas las funciones reales (o complejas) que son continuas en el intervalo cerrado $[0, 1]$.
        $$C^0[0,1] = \{ f: [0,1] \to \mathbb{R} \mid f \text{ es continua en } [0,1] \}$$
    \end{observacion}
    
    Este operador es lineal y continuo. Para verificar la acotación:
    \begin{equation*}
        \|T(f)\|_\infty = \sup_{t \in [0,1]} \left| \int_{0}^{t} f(x) \, dx \right| \leq \sup_{t \in [0,1]} \int_{0}^{t} |f(x)| \, dx \overset{(*)}{\leq} \int_{0}^{1} |f(x)| \, dx \leq \|f\|_\infty
    \end{equation*}
    donde en $(*)$ hemos usado que $|f(x)|\geq 0$ para todo $x \in [0,1]$, y por tanto es no negativo. Dado que $\|T(f)\|_\infty \leq \|f\|_\infty$, el operador es acotado con norma $\|T\| \leq 1$ y, por tanto, continuo. \\

    \noindent
    Más en general, sea $T: C^0([a,b]) \longrightarrow C^0([a,b])$ definido mediante un núcleo $K(x,t)$ continuo en $[a,b] \times [a,b]$:
    \begin{equation*}
        T(f)(x) = \int_{a}^{b} K(x,t)f(t) \, dt
    \end{equation*}
    Usando la desigualdad de Minkowski (o propiedades básicas de la integral):
    \begin{align*}
        |T(f)(x)| \leq& \int_{a}^{b} |K(x,t)| \cdot |f(t)| \, dt \overset{(*)}{\leq}\int_{a}^{b} |K(x,t)| \cdot \|f\|_\infty \, dt \overset{(**)}{\leq} \|f\|_\infty \int_{a}^{b} \|K\|_\infty \, dt = \\
        =& \|f\|_\infty \cdot \|K\|_\infty \int_{a}^{b} dt = |b-a| \cdot \|K\|_\infty \cdot \|f\|_\infty
    \end{align*}
    donde en $(*)$ hemos acotado a $|f(t)|$ por $\|f\|_\infty= \sup_{t \in [a,b]} |f(t)|$, es decir, su máximo valor absoluto en el intervalo $[a,b]$, y en $(**)$ hemos hecho lo mismo para $|K(x,t)|$.  \\
    \noindent
    Definiendo $M = |b-a| \cdot \|K\|_\infty$, se tiene $\|Tf\|_\infty \leq M \|f\|_\infty$, lo que demuestra que $T$ es continuo.
    
    \item Sea $F_k: l_p \longrightarrow \mathbb{R}$ el funcional que asigna a cada sucesión su $k$-ésima componente, es decir, $F_k(x) := x_k$. Este operador es continuo, ya que:
    \begin{equation*}
        \|F_k(x)\| = |x_k| = \left( |x_k|^p \right)^{1/p} \leq \left( \sum_{l=1}^{\infty} |x_l|^p \right)^{1/p} = \|x\|_{p}
    \end{equation*}
    Esto muestra que $\|F_k\| \leq 1$, por lo que el operador está acotado y es, por tanto, continuo.
\end{enumerate}
\end{ejemplo}

\begin{definicion}
    Sean $X, Y$ espacios normados. Definimos el \tbf{conjunto de los operadores lineales acotados} como:
    \begin{equation*}
        BL(X, Y) := \{ T: X \longrightarrow Y \mid T \text{ es lineal y acotado} \}
    \end{equation*}
\end{definicion}

\begin{prop}\label{prop:norma-operacional}
    $BL(X, Y)$ es un espacio vectorial normado con la norma operacional definida por:
    \begin{equation*}
        \|T\| = \sup_{x \neq 0} \frac{\|T(x)\|}{\|x\|} = \sup_{\|x\|=1} \|T(x)\|, \quad \forall T \in BL(X, Y)
    \end{equation*}

\begin{proof}
    Observamos primero que si $T$ es acotado, existe $M > 0$ tal que:
    \begin{equation*}
        \|Tx\| \leq M\|x\| \implies \frac{\|Tx\|}{\|x\|} \leq M \quad \forall x \neq 0 \implies \|T\| \leq M < +\infty
    \end{equation*}
    Para demostrar que $\|T\|$ es una norma, verificamos la desigualdad triangular, pues los otros dos axiomas, $\|\lambda T\| = |\lambda|\|T\|$ y $\|T\|=0 \iff T=0$, se cumplen de forma trivial, veámoslo:

    \begin{itemize}
        \item $\|\lambda T\| = \sup\limits_{\|x\|=1} \|\lambda (Tx)\|=\sup\limits_{\|x\|=1} \left(|\lambda| \|Tx\|\right)= |\lambda| \sup\limits_{\|x\|=1} \|Tx\| = |\lambda|\|T\|$ para todo $\lambda \in \mathbb{R}$ y $T \in BL(X, Y)$.
        \item $\|T\|=0 \iff \sup\limits_{\|x\|=1} \|T(x)\|=0 \iff T(x)=0 \text{ para todo } x \text{ con } \|x\|=1 \iff T=0$.
    \end{itemize}

    \noindent
    Visto los axiomas sencillo, veámos finalmente la desigualdad triangular. Sean $T_1, T_2 \in BL(X, Y)$. Para cualquier $x \in X$, tenemos:
    \begin{align*}
        \|(T_1 + T_2)(x)\| &= \|T_1x + T_2x\| \\
        &\leq \|T_1x\| + \|T_2x\| \\
        &\leq \|T_1\| \cdot \|x\| + \|T_2\| \cdot \|x\| \\
        &= (\|T_1\| + \|T_2\|) \cdot \|x\|
    \end{align*}
    Tomando $\|x\|=1$, obtenemos:
    \begin{equation*}
        \|(T_1 + T_2)(x)\| \leq \|T_1\| + \|T_2\|, \quad \forall x\in X
    \end{equation*}
    donde sabiendo que $\| T_1 + T_2\| = \sup\limits_{\|x\|=1} \|(T_1 + T_2)(x)\|$, y tomando que la desigualdad anterior se cumple para todo $x\in X$ tal que $\|x\|=1$, entonces se tiene que $\|T_1 + T_2\| \leq \|T_1\| + \|T_2\|$, lo que demuestra la desigualdad triangular.
\end{proof}
\end{prop}

\begin{observacion}
    Realmente la norma $\|T\|$ es la constante más pequeña $M$ tal que $\|Tx\| \leq M\|x\|$.
\end{observacion}

\begin{notacion}
    Importante notar, que a veces en lugar de poner $\|T(x)\|$ se suele escribir $\|Tx\|$, es decir, se omite el paréntesis, y se escribe $Tx$ en lugar de $T(x)$, aunque ambos significados son equivalentes.
\end{notacion}

\observacion{
    Para ver más claramente la última igualdad de la Proposición \ref{prop:norma-operacional}, es decir, la siguiente
    \begin{equation*}
        \sup_{x \neq 0} \frac{\|T(x)\|}{\|x\|} = \sup_{\|x\|=1} \|T(x)\|
    \end{equation*}
    basta observar el siguiente hecho, tomando $u = \frac{x}{\|x\|}$, entonces 
    \begin{equation*}
        \|u\| = \left\| \frac{x}{\|x\|} \right\| = \frac{1}{\|x\|} \cdot \|x\| = 1
    \end{equation*}
    Aplicando esto, tenemos que 
    $$\frac{\|T(x)\|}{\|x\|} = \frac{1}{\|x\|} \cdot \|T(x)\| \overset{(*)}{=} \left\| T \left( \frac{1}{\|x\|} \cdot x \right) \right\|$$
    donde en $(*)$ hemos usado la propiedad de escalares de la norma (aunque no hemos probado que esta sea una norma ya, si hemos probado dicha propiedad justo antes).
}

\begin{definicion}
    Como $BL(X, Y)$ es normado, entonces tenemos la noción de convergencia. $(T_n) \in BL(X, Y)$ \tbf{converge (o converge uniformemente)} a $T \in BL(X, Y)$ si:
    \begin{equation*}
        \|T_n - T\| \to 0 \iff \sup_{\|x\|\neq 0} \dfrac{\|(T_n - T)(x)\|}{\|x\|} \to 0
    \end{equation*}
\end{definicion}

\observacion{
    Si $T_n \to T$ entonces $T_n(x) \to T(x) \, \forall x \in X$, de hecho:
    \begin{equation*}
        ||T_n(x) - T(x)|| = ||(T_n - T)(x)|| \le ||T_n - T|| \cdot ||x|| \to 0
    \end{equation*}
}

\begin{definicion}
    Llamamos $BL(X, \mathbb{R}) =: X^*$ al \textbf{espacio dual continuo} de $X$. \\
\end{definicion}

\begin{teo}
    Sea $Y$ un espacio de Banach, y $X$ un espacio normado. Entonces $BL(X, Y)$ es de Banach.
\begin{proof}
    Sea $(T_n)_{n\in \N}$ una sucesión de Cauchy de $BL(X, Y)$, es decir: $||T_n - T_m|| \le \varepsilon \ \forall n, m > N$. Entonces:
    
    \begin{equation*}
        \forall x \in X \quad \quad \|T_n x - T_m x\| \le ||T_n - T_m|| \cdot ||x|| \le \varepsilon \cdot ||x|| \ \forall n, m > N
    \end{equation*}
    Como $\{T_n(x)\}_{n\in \N}$ es de Cauchy en $Y$, que es completo, entonces $\exists T(x)$ tal que:
    \begin{equation*}
        T_n(x) \to T(x) \text{ cuando } n \to +\infty
    \end{equation*}
    Sea $T$ la aplicación $x \mapsto T(x)$, entonces: 
    \begin{enumerate}
        \item[(i)] $T$ es lineal:
        \begin{equation*}
            T(x+y) = \lim_{n \to \infty} T_n(x+y) = \lim_{n \to \infty} (T_n(x) + T_n(y)) = T(x) + T(y)
        \end{equation*}
        \item[(ii)] $T$ es acotado: \\
        Dado que $\{T_n(x)\}$ es de Cauchy en $Y$, entonces $\{||T_n(x)||\} \subset \mathbb{R}$ es una sucesión de números reales que es de Cauchy (esto ocurre porque la norma es una función continua y mantiene la sucesión de Cauchy), y por lo tanto es acotada, es decir, $\exists M > 0 : ||T_n(x)|| \le M$. Entonces:
        \begin{equation*}
            ||T_n(x)|| \le ||T_n|| \cdot ||x|| \le M \cdot ||x|| \, \forall x \in X
        \end{equation*}
        Pasando al límite:
        \begin{equation*}
            ||T(x)|| \le M \cdot ||x|| \Rightarrow T \in BL(X, Y)
        \end{equation*}
        \item[(iii)] $T_n \to T$ en $BL(X, Y)$: \\
        Ya sabemos que $||T_n(x) - T_m (x)|| \le \varepsilon \cdot ||x|| \ \forall n, m > N$. Hacemos $n \to +\infty$:
        \begin{equation*}
            ||T(x) - T_m(x)|| \le \varepsilon \cdot ||x|| \iff \frac{||T(x) - T_m(x)||}{||x||} \le \varepsilon \quad \text{ con } \|x\|\neq 0
        \end{equation*}
        donde hemos supuesto $\|x\|\neq 0$ porque si $\|x\|=0$ entonces se cumpliría la desigualdad trivialmente. Se tiene finalmente que $T_m \to T$.
    \end{enumerate}
\end{proof}
\end{teo}

\begin{coro}
    $X^*$ es de Banach para todo espacio normado $X$, sea o no de Banach.
\end{coro}

\noindent
Estudiaremos ahora el caso $Y=X$, donde primero definiremos $BL(X) := BL(X, X)$, y luego veremos que $BL(X)$ es cerrado para la composición. \\
\noindent
Sean $T_1, T_2 \in BL(X)$, tales que $X \xrightarrow{T_1} X \xrightarrow{T_2} X$. Entonces:
\begin{equation*}
    T_2 \circ T_1: X \longrightarrow X, \quad (T_2 \circ T_1)(x) = T_2(T_1(x)) \quad \forall x \in X
\end{equation*}
Veámos ahora si está acotada esta composición 
\begin{equation*}
    \dfrac{\|(T_2 \circ T_1) (x)\|}{\|x\|} = \frac{\|T_2(T_1 x)\|}{\|T_1 x\|} \cdot \frac{\|T_1 x\|}{\|x\|} \le \|T_2\| \cdot \|T_1\|
\end{equation*}
donde hemos supuesto $\|x\|\neq 0$ y $\|T_1 x\|\neq 0$, y se ha obtenido que $\|T_2 \circ T_1\| \le \|T_2\| \cdot \|T_1\| < +\infty$, lo que demuestra que $T_2 \circ T_1$ es acotado, y por tanto, $T_2 \circ T_1 \in BL(X)$. \\ \\

\noindent
Comprobaremos ahora que la composición también es \tbf{continua}, es decir, el límite de la composición es la composición de los límites. \\
\noindent
Sean $(T_{1,n}), (T_{2,n}) \in BL(X)$ tales que $T_{1,n} \to T_1$ y $T_{2,n} \to T_2$, entonces:
\begin{align*}
\|T_{2,n} \circ T_{1,n} - T_2 \circ T_1\| &= \|T_{2,n} \circ T_{1,n} - T_{2,n} \circ T_1 + T_{2,n} \circ T_1 - T_2 \circ T_1\| \\
&\le \|T_{2,n}(T_{1,n} - T_1)\| + \|(T_{2,n} - T_2) \circ T_1\| \\
&\le \|T_{2,n}\| \cdot \|T_{1,n} - T_1\| + \|T_{2,n} - T_2\| \cdot \|T_1\| \to 0
\end{align*}
por lo que $T_{2,n} \circ T_{1,n} \to T_2 \circ T_1$. Esto demuestra que la composición es continua, lo que implica que $(BL(X), \circ)$ es un \textbf{Álgebra de Banach} con elemento neutro $id_X$.

\begin{definicion}
    Para que un conjunto $\mathcal{A}$ sea un \tbf{Álgebra de Banach}, debe cumplir cuatro condiciones fundamentales:
    \begin{enumerate}
        \item Es un espacio vectorial
        \item Existe una operación de ``multiplicación" entre sus elementos, que es 
        \begin{itemize}
            \item \tbf{Asociativa}: $(a \cdot b) \cdot c = a \cdot (b \cdot c)\quad \forall a, b, c \in \mathcal{A}$.
            \item \textbf{Distributiva}: $a \cdot (b + c) = a \cdot b + a \cdot c \quad \forall a, b, c \in \mathcal{A}$.
        \end{itemize}
        \item Es un espacio de Banach.
        \item Para cualquier par de elementos $a, b \in \mathcal{A}$, se debe cumplir:$$\|ab\| \leq \|a\| \cdot \|b\|$$
    \end{enumerate}
\end{definicion}

\begin{ejercicio}
    Sea $F_k(x) = x_k$ con $x \in \ell_p$. Mostrar que $\{F_k x\}$ converge puntualmente a $0$, pero $\|F_k\| \not\to 0$. \\
    \noindent
    Como $x \in \ell_p$, sabemos que $\lim\limits_{k \to \infty} |x_k|^p = 0$, y tomando $f(y)= y^p$ función continua en $y\geq 0$, tenemos que 
    $$\lim_{k \to \infty} |x_k| = \lim_{k \to \infty} \left( |x_k|^p \right)^{1/p} = \left( \lim_{k \to \infty} |x_k|^p \right)^{1/p} = 0^{1/p} = 0$$
    y por tanto 
    $$\lim_{k \to \infty} |x_k| = 0 \implies \lim_{k \to \infty} x_k = 0$$
    Tenemos ahora la convergecia puntual de la sucesión $\{F_k x\}$ a $0$
    \begin{equation*}
        \lim_{k \to \infty} F_k x = \lim_{k \to \infty} x_k = 0 \quad (\text{Convergencia puntual})
    \end{equation*}
    Por otro lado, observamos que:
    \begin{equation*}
        |F_k(x)| = |x_k|= (|x_k|^p)^{1/p}\le \left( \sum_{i=1}^{\infty} |x_i|^p \right)^{1/p} = \|x\|_{p} \implies |F_k(x)| \le 1 \cdot \|x\|_{p} \quad \forall x \in \ell_p   
    \end{equation*}
    de lo que tenemos $\|F_k\| \le 1$. Si tomamos la sucesión $x = e_k = (0, 0, \dots, 1, 0, \dots)$ (donde el 1 está en la posición $k$):
    \[
    \|e_k\|_{\ell_p} = 1 \quad \text{y} \quad F_k(e_k) = 1 \implies \|F_k\| \ge \frac{|F_k(e_k)|}{\|e_k\|} = 1
    \]
    De lo anterior concluimos que $\|F_k\| = 1$ para todo $k$, por lo que la sucesión de operadores \textbf{no} converge a 0 en norma, aunque sí lo haga puntualmente.
\end{ejercicio}

\begin{ejercicio}
    Sea $F: C^0([a,b]) \to \mathbb{R}$ definido por $F(u) = u(c)$ para un $c \in [a,b]$ fijo. Calcular $\|F\|$. \\
    \noindent
    Sea $u \in C^0([a,b])$. Como $[a,b]$ es compacto, $u$ alcanza su máximo, por lo que:
    \[
    |F(u)| = |u(c)| \le \max_{t \in [a,b]} |u(t)| = \|u\|_\infty \implies \|F\| \le 1
    \]
    Si tomamos la función constante $u(t) = 1$, tenemos $\|u\|_\infty = 1$ y $|F(1)| = 1$. Por lo tanto, $\|F\| = 1$, donde hemos usado que 
    \begin{equation*}
        \|F\| = \sup_{\|u\|_\infty = 1} |F(u)| \leq 1
    \end{equation*}
    y por tanto si encontramos una función $u\in C^0([a,b])$ que cumpla $F(u)=1$, tenemos el supremo buscado, es decir, $\|F\|=1$.
\end{ejercicio}

\begin{ejercicio}
\begin{itemize}
    \item[(i)] Sea $F: C^0([a,b]) \to \mathbb{R}$ tal que $F(u) = \int_{a}^{b} u(t) \, dt$. Calcular $\|F\|$.\\
    Calculemos la norma
    \begin{equation*}
        |F(u)| = \left| \int_{a}^{b} u(t) \, dt \right| \le \int_{a}^{b} |u(t)| \, dt \le \|u\|_\infty \int_{a}^{b} 1 \, dt = |b-a| \cdot \|u\|_\infty
    \end{equation*}
    Esto implica $\|F\| \le |b-a|$. Si $u(t) = 1$, entonces $F(1) = b-a$, luego $\|F\| = |b-a|$.

    \item[(ii)] Sea $T: C^0([a,b]) \to C^0([a,b])$ tal que $T(u)(t) = \int_{a}^{t} u(x) \, dx$. Calcular $\|T\|$.\\
    Calculemos la norma del operador:
    \begin{equation*}
        \|T(u)\|_\infty = \sup_{t \in [a,b]} \left| \int_{a}^{t} u(x) \, dx \right| \le \int_{a}^{b} |u(x)| \, dx \le |b-a| \cdot \|u\|_\infty \implies \|T\| \le |b-a|
    \end{equation*}
    Si $u(t) = 1$, entonces $T(1)(t) = t-a$. El máximo se alcanza en $t=b$:
    \begin{equation*}
        \|T(1)\|_\infty = \sup_{t \in [a,b]} |t-a| = |b-a| \implies \|T\| = |b-a|
    \end{equation*}
\end{itemize}
\end{ejercicio}

\begin{ejercicio}
    Sea $T: \ell_2 \to \ell_2$ el operador de desplazamiento a la izquierda (Shift) definido por $T_k(x) = (x_{k+1}, x_{k+2}, \dots)$. Demostrar que $T_k$ es un operador lineal acotado y calcular su norma. \\
    \noindent
    Demostremos que es acotado
    \begin{equation*}
        \|T_k(x)\|_2^2 = \left(\sqrt{\sum_{j=k+1}^{\infty} |x_j|^2} \right)^2 \le \sum_{j=1}^{\infty} |x_j|^2 = \|x\|_2^2 \implies \|T_k(x)\|_2 \le \|x\|_2
    \end{equation*}
    Esto muestra que $\|T_k\| \le 1$. Para ver que la norma es exactamente 1, tomamos el vector de la base canónica $e_{k+1} = (0, \dots, 0, 1, 0, \dots)$ (donde el 1 está en la posición $k+1$):
    \begin{equation*}
        \|e_{k+1}\|_2 = 1 \quad \text{y} \quad T_k(e_{k+1}) = (1, 0, 0, \dots) \implies \|T_k(e_{k+1})\|_2 = 1
    \end{equation*}
    Por lo tanto, $\|T_k\| = 1$.
\end{ejercicio}

\subsection{Distancias equivalentes}

\begin{definicion}
    Sea $X$ un espacio vectorial. Dos espacios normados $(X, \|\cdot\|_1)$ y $(X, \|\cdot\|_2)$ se dicen \textbf{equivalentes} si son equivalentes como espacios métricos (\tbf{equivalencia métrica}), es decir, si y solo si cada conjunto abierto en $(X, \|\cdot\|_1)$ contiene un conjunto abierto de $(X, \|\cdot\|_2)$ y viceversa (\tbf{equivalencia topológica}).
\end{definicion}

\begin{observacion}
    En términos generales, el hecho de que dos espacios sean equivalentes no garantiza que preserven el carácter de las sucesiones de Cauchy (una sucesión de Cauchy en un espacio podría no serlo en el otro).
\end{observacion}

\begin{ejemplo}
    Consideremos $(\mathbb{R}, d_1)$ con la métrica usual $d_1(x, y) = |x - y|$ y $(\mathbb{R}, d')$ con la métrica $d'(x, y) = \arctan(|x - y|)$. Aunque ambos espacios son topológicamente equivalentes, observamos que:
    \begin{itemize}
        \item La sucesión $x_n = n$ es de Cauchy en $(\mathbb{R}, d')$ ya que $\|x_n\|_{d'} = \arctan n \longrightarrow \frac{\pi}{2}$, lo que implica su convergencia en la completación del espacio.
        \item Sin embargo, dicha sucesión no converge (ni es de Cauchy) en $(\mathbb{R}, d_1)$.
    \end{itemize}
\end{ejemplo}

\begin{definicion}
    Sean $(X, d_1)$ y $(X, d_2)$ dos espacios métricos. Se dice que las métricas $d_1$ y $d_2$ son \textbf{uniformemente equivalentes} (o simplemente \tbf{equivalentes}) si
    \begin{equation*}
        \exists \alpha, \beta > 0 \text{ tales que } \alpha d_1(x, y) \leq d_2(x, y) \leq \beta d_1(x, y) \quad \forall x, y \in X
    \end{equation*}
\end{definicion}

\notacion{
    Es importante destacar que cuando hablemos de ``uniformemente equivalentes'' nos referiremos a la \tbf{equivalencia métrica} (definición anterior), mientras que si hablamos de ``equivalentes'' nos referiremos a la \tbf{equivalencia topológica}.
}

\begin{prop}
    Si $(X, d_1)$ y $(X, d_2)$ son uniformemente equivalentes, entonces:
    \begin{equation*}
        (x_n) \text{ es de Cauchy en } d_1 \iff (x_n) \text{ es de Cauchy en } d_2
    \end{equation*}
\end{prop}

\begin{proof}
    Supongamos que $(x_n)$ es una sucesión de Cauchy en $(X, d_1)$. Por definición, para todo $\varepsilon > 0$ existe $N$ tal que $d_1(x_n, x_m) \leq \varepsilon$ para todo $n, m > N$. Utilizando la equivalencia uniforme:
    \begin{equation*}
        d_2(x_n, x_m) \leq \beta d_1(x_n, x_m) \leq \beta \varepsilon
    \end{equation*}
    Dado que $\beta \varepsilon$ puede hacerse arbitrariamente pequeño, se deduce que $(x_n)$ es de Cauchy en $d_2$. El recíproco se demuestra de forma análoga utilizando la constante $\alpha$.
\end{proof}

\noindent
En general, si dos distancias son equivalentes, no tienen porque ser uniformemente equivalentes, cuando hablamos de espacios métricos, mientras que en espacios normados si es lo mismo, lo veremos en la siguiente proposición.

\begin{prop}
    Dos normas $\|\cdot\|_1$ y $\|\cdot\|_2$ sobre un espacio vectorial $X$ son equivalentes si y solo si existen constantes $\alpha, \beta > 0$ tales que:
    \begin{equation*}
        \alpha \|x\|_1 \leq \|x\|_2 \leq \beta \|x\|_1 \quad \forall x \in X
    \end{equation*}

\begin{proof}
    \begin{description}
        \item[$\Longleftarrow)$] Si existen tales constantes, las métricas inducidas $d_1(x,y) = \|x-y\|_1$ y $d_2(x,y) = \|x-y\|_2$ son uniformemente equivalentes. Esto garantiza que ambas normas definen la misma topología (mismos conjuntos abiertos) y, por tanto, son equivalentes.
        \item[$\Longrightarrow)$] Supongamos que las normas son equivalentes. Esto implica que la aplicación identidad:
        \begin{equation*}
            Id: (X, \|\cdot\|_1) \longrightarrow (X, \|\cdot\|_2)
        \end{equation*}
        es continua. Esta conclusión viene de una noción topológica, donde se tiene que una aplicación es continua, si tomando un abierto en $(X, \|\cdot\|_2)$, su imagen inversa es un abierto en $(X, \|\cdot\|_1)$. Dado que las normas son equivalentes, un conjunto abierto $V$ en el espacio de llegada, en el espacio inicial es $Id^{-1}(V)=V$, que es un conjunto abierto para la norma $\|\cdot\|_1$, pues dijimos que las normas eran equivalentes. \\

        \noindent
        Dado que el operador identidad es claramente lineal, la continuidad implica que es un operador acotado. Por lo tanto, existe una constante $\beta > 0$ tal que:
        \begin{equation*}
            \|Id(x)\|_2 \leq \beta \|x\|_1 \implies \|x\|_2 \leq \beta \|x\|_1 \quad \forall x \in X
        \end{equation*}
        Por simetría, la aplicación inversa $Id^{-1}: (X, \|\cdot\|_2) \longrightarrow (X, \|\cdot\|_1)$ también es lineal y continua, lo que implica que es acotada. Así, existe una constante $\gamma > 0$ tal que:
        \begin{equation*}
            \|Id^{-1}(x)\|_1 \leq \gamma \|x\|_2 \implies \|x\|_1 \leq \gamma \|x\|_2 \quad \forall x \in X
        \end{equation*}
        De esta última desigualdad obtenemos $\frac{1}{\gamma} \|x\|_1 \leq \|x\|_2$. Definiendo $\alpha = \frac{1}{\gamma}$, se llega a la relación:
        \begin{equation*}
            \alpha \|x\|_1 \leq \|x\|_2 \leq \beta \|x\|_1
        \end{equation*}
    \end{description}
\end{proof}    
\end{prop}

\observacion{
    La equivalencia entre \underline{espacios normados} puede ser vista como equivalencia métrica o topológica, en ambos casos es equivalente como hemos comprobado en la anterior proposición. En cambio en \underline{espacios métricos}, la equivalencia topológica es la que habla de abiertos y la métrica es la que habla de desigualdades entre las métricas.
}


\subsection{Convergencia de series de operadores}

\noindent 
Sea $X$ un espacio de Banach y sea $(T_n)_{n\in\N} \in BL(X)$ una sucesión. Definimos 
$$S_n := \sum_{k=1}^{n} T_k$$

\begin{definicion}
    Se dice que la sucesión $(S_n)_{n\in\N}$ es convergente si exsite el límite
    \begin{equation*}
        \lim_{n \to +\infty} S_n =: T
    \end{equation*}
    es decir, si $\|S_n - T\|_{BL(X)} \xrightarrow{n \to +\infty} 0$.
\end{definicion}

\begin{definicion}
    Una serie de operadores $(S_n)_{n\in\N}$ se dice \tbf{absolutamente convergente} si 
    $$\sum_{n=1}^{\infty} \|T_n\|$$
    es convergente como sucesión de números reales.
\end{definicion}

\observacion{
    Si $(S_n)$ converge absolutamente, entonces converge.
    \begin{proof}
        \begin{equation*}
            \left\|\sum_{k=1}^{n} T_k \right\| \leq \sum_{k=1}^{n} \|T_k\| \quad \text{de lo cual} \quad \left\|\sum_{k=1}^{+\infty} T_k\right\| < +\infty \implies S_n \text{ es convergente}
        \end{equation*}
    \end{proof}
}

\begin{ejercicio}
    Sea $T\in BL(X)$ con $\|T\| < 1$. Demostrar que 
    $$\sum_{k=0}^{+\infty} T^k=(Id-T)^{-1}$$
    \noindent
    Primero demostraremos que converge absolutamente, donde sabemos que $(BL(X), \|\cdot\|)$ es un Algebra de Banach, entonces tenemos que la norma de operadores tiene la propiedad de submultiplicatividad, es decir, $\|T^k\| \leq \|T\|^k$, entonces 
    \begin{equation*}
        \sum_{k=0}^{+\infty} \|T^k\| \leq \sum_{k=0}^{+\infty} \|T\|^k
    \end{equation*}
    y como $\|T\| < 1$, entonces $\sum\limits_{k=0}^{+\infty} \|T\|^k$ es una serie geométrica convergente, lo que implica que $\sum\limits_{k=0}^{+\infty} \|T^k\|$ es convergente, lo que implica que $\sum\limits_{k=0}^{+\infty} T^k$ es convergente. \\
    \noindent
    Demostrado que la sucesión es convergente (existe el límite), calcularemos cual es el límite, veámoslo
    \begin{align*}
        \sum_{k=0}^{+\infty} T^k &= Id + \sum_{k=1}^{+\infty} T^k = Id + T \cdot \sum_{k=1}^{+\infty} T^{k-1} = Id + T \sum_{k=0}^{+\infty} T^k \\
        &\implies (Id - T) \cdot \sum_{k=0}^{+\infty} T^k = Id \implies \sum_{k=0}^{+\infty} T^k = (Id - T)^{-1}
    \end{align*}
\end{ejercicio}


\begin{ejercicio}
    Sea $T: X \to Y$ lineal. Demostrar que 
    \begin{equation*}
        T \text{ es continua} \iff T \text{ manda sucesiones de Cauchy a sucesiones de Cauchy}
    \end{equation*}
    \begin{description}
        \item[$\Longrightarrow)$] Sea $(x_n)$ una sucesión de Cauchy. Entonces $\|x_n - x_m\| \xrightarrow{n, m \to +\infty} 0$. De lo cual
        \begin{equation*}
            \|T(x_n) - T(x_m)\| = \|T(x_n - x_m)\| \leq \|T\| \cdot \|x_n - x_m\| \xrightarrow{n,m \to +\infty} 0
        \end{equation*}
        donde hemos usado que como $T$ es un operador lineal continuo, entonces acotado y por tanto $\|T\| < +\infty$. Queda demostrado que $(T(x_n))_n$ es de Cauchy. \\
        \item[$\Longleftarrow)$] Supongamos por reducción al absurdo que $T$ no es continua, entonces $T$ no es acotada (lo cual es equivalente a no ser continua por sucesiones), es decir, 
        \begin{equation*}
            \text{ para cualquier constante } M > 0, \text{ existe algún vector } v \in X \text{ tal que } \|T(v)\| > M\|v\|
        \end{equation*}
        Sabiendo esto, obtendremos la sucesión $(z_n)$ tal que 
        \begin{equation*}
            \|T(z_n)\| > n^2 2^n \|z_n\| 
        \end{equation*}
        y definirimos la sucesión $(x_n)$ de la siguiente manera
        \begin{equation*}
            x_n =\frac{1}{n^2}\cdot \frac{z_n}{\|z_n\|}
        \end{equation*}
        Veámos que dos desigualdades obtenemos de la definición de $x_n$, tanto para la norma de los elementos de la sucesión, como para la norma de su imagen por $T$
        \begin{itemize}
            \item $\|x_n\| = \left\| \dfrac{1}{n^2} \dfrac{z_n}{\|z_n\|} \right\| = \dfrac{1}{n^2} \cdot \dfrac{\|z_n\|}{\|z_n\|} = \dfrac{1}{n^2}$
            \item $\|T(x_n)\| = \left\| T\left( \dfrac{1}{n^2} \cdot \dfrac{z_n}{\|z_n\|} \right) \right\| = \dfrac{1}{n^2} \cdot \dfrac{\|T(z_n)\|}{\|z_n\|} > \dfrac{1}{n^2} \cdot \dfrac{n^2 2^n \|z_n\|}{\|z_n\|} = 2^n$
        \end{itemize}
        Entonces, sea $S_n = \sum\limits_{k=1}^{n} x_k$ tal que 
        $$\|S_n -S_m\|= \left\| \sum\limits_{k=m+1}^{n} x_k \right\|\leq \sum\limits_{k=m+1}^{n} \|x_k\| = \sum\limits_{k=m+1}^{n} \dfrac{1}{k^2} \xrightarrow[n,m \to +\infty]{} 0$$
        es de Cauchy. Por otro lado, 
        $$\|T(S_n)\| = \|\sum_{k=1}^{n} T(x_k)\| > \sum_{k=1}^{n} 2^k \to +\infty$$
        lo que implica que $(T(S_n))_n$ no es de Cauchy, ya que si la sucesión no está acotada, entonces no es de Cauchy. Esto es una contradicción, por lo tanto, $T$ es continua. \\
    \end{description}
\end{ejercicio}


\section{Espacios normados de dimensión finita}

\begin{definicion}
    Un espacio normado $(X, \|\cdot\|)$ se dice que es de \textbf{dimensión finita} si el espacio vectorial subyacente $X$ es finito, es decir, posee una base formada por un número finito de elementos.
\end{definicion}

\begin{prop}(Desigualdad de Cauchy-Schwarz)
    Sean $a_1, \dots, a_n$ y $b_1, \dots, b_n$ números reales cualesquiera. Entonces se cumple la siguiente desigualdad:
    \begin{equation*}
        \left( \sum_{k=1}^n a_k b_k \right)^2 \leq \left( \sum_{k=1}^n a_k^2 \right) \left( \sum_{k=1}^n b_k^2 \right)
    \end{equation*}
    Además, la igualdad se verifica si y sólo si existe un número real $x$ tal que
    \begin{equation*}
        a_k x + b_k = 0 \quad \text{ para cada } k = 1, \dots, n.
    \end{equation*}
\end{prop}

\begin{prop}
    Sea $(X, \|\cdot\|)$ un espacio normado con $\dim X = N < +\infty$ y base $\{\alpha_1, \dots, \alpha_N\}$. Entonces la aplicación:
    \begin{equation*}
        T : X \longrightarrow \mathbb{R}^N, \quad x = \sum_{i=1}^N a_i \alpha_i \mapsto 
        \begin{pmatrix} 
            a_1 \\ 
            \vdots \\ 
            a_N 
        \end{pmatrix}
    \end{equation*}
    es un isomorfismo lineal topológico.
\begin{proof}
    $T$ es un isomorfismo lineal por ser $X$ un espacio vectorial de dimensión finita, esto se ve fácilmente demostrando las propiedades siguientes:
    \begin{itemize}
        \item \tbf{Linealidad:} $T(ax + by) = aT(x) + bT(y)$. Esto es directo por la definición de coordenadas.
        \item \tbf{Bijectividad:} Es \underline{inyectiva} porque en una base la representación es única (si las coordenadas son cero, el vector es cero). Es \underline{sobreyectiva} porque para cualquier columna de números en $\mathbb{R}^N$, puedes construir un vector en $X$ sumando la base.
    \end{itemize}
    Para demostrar que $T$ es un isomorfismo topológico, tenemos que demostrar que $T$ y su inversa $T^{-1}$ son continuas, para ello usaremos la Proposición \ref{prop:continuidad-acotado}, ya que sabemos que $T$ y $T^{-1}$ son operadores lineales (se puede demostrar trivialmente sabiendo que trabajamos con una base).
    \begin{itemize}        
        \item \textbf{Continuidad de $T^{-1}$:} Consideramos la aplicación inversa dada como 
        $$T^{-1} : \mathbb{R}^N \longrightarrow X, \quad T^{-1}(y) = \sum_{i=1}^N a_i \alpha_i\quad \forall y = (a_1, \dots, a_N)^T\in \R^N$$
        
        Usando la desigualdad de Cauchy-Schwarz:
        \begin{equation*}
            \|T^{-1}(y)\| = \left\| \sum_{i=1}^N a_i \alpha_i \right\| \leq \sum_{i=1}^N |a_i| \|\alpha_i\|\leq \left( \sum_{i=1}^N |a_i|^2 \right)^{1/2} \left( \sum_{i=1}^N \|\alpha_i\|^2 \right)^{1/2}= \|y\| M
        \end{equation*}
        Si definimos $M = (\sum\limits_{i=1}^N \|x_i\|^2)^{1/2}$ como una constante (los vectores de la base no cambian), tenemos $\|T^{-1}(y)\| \leq M \|y\|_2$. Al ser un operador acotado, $T^{-1}$ es continuo.

        \item \textbf{Continuidad de $T$:} Supongamos por reducción al absurdo que $T$ no es continua en el origen. Entonces buscamos una sucesión $(x_n) \subset X$ tal que $x_n \to 0$ pero $\|T(x_n)\| \not\to 0$, y que cumpla que

        \begin{equation*}
                \exists \varepsilon > 0 \text{ tal que } \|T(x_n)\| > \varepsilon \text{ para todo } n
        \end{equation*}
        Para ello sabemos que $T$ no es continua en el origen, es decir, existe un $\bar{\varepsilon} > 0$ tal que para cualquier $\delta > 0$ que yo elija, siempre hay al menos un vector $x$ que cumple $\|x\| < \delta$ y $\|T(x)\| > \bar{\varepsilon}$, y como es cierto para cualquier $\delta$, construiremos la sucesión de la que hablabamos antes de la siguiente manera

        \begin{itemize}
            \item Para $\delta = 1$, existe $x_1$ tal que $\|x_1\| < 1$ y $\|T(x_1)\| > \bar{\varepsilon}$.
            \item Para $\delta = \frac{1}{2}$, existe $x_2$ tal que $\|x_2\| < \frac{1}{2}$ y $\|T(x_2)\| > \bar{\varepsilon}$.
            \item Y así sucesivamente para cualquier $n$.
        \end{itemize}
        Finalmente tenemos que una sucesión $\{x_n\}_{n\in \N}$ tal que $x_n \to 0$, pues $\|x_n\| \leq \frac{1}{n}$ pero $\|T(x_n)\| > \bar{\varepsilon}$ para todo $n$. \\
        
        \noindent
        Definimos $z_n = \dfrac{x_n}{\|T(x_n)\|}$, que lo podemos hacer porque sabemos que el denominador nunca es $0$ (por lo demostrado antes). Se observa que $\|T(z_n)\| = 1$ para todo $n$, y que $z_n\to 0$ porque 
        \begin{equation*}
            \|z_n\| = \frac{\|x_n\|}{\|T(x_n)\|} < \frac{1/n}{\bar{\varepsilon}} \to 0
        \end{equation*} 
        
        \noindent
        La sucesión $\{T(z_n)\}$ está contenida en la esfera unidad $S^{N-1}=\left\{x\in \mathbb {R} ^{n}:\left\|x\right\|=1\right\}$ de $\mathbb{R}^N$, que es un conjunto compacto. Por tanto, existe una subsucesión $\{T(z_{n_k})\}$ que converge a un elemento $y \in \mathbb{R}^N$ con $\|y\| = 1$, esto es cierto gracias la continuidad de $T$, es decir

        \begin{equation*}
            \|y\| = \lim_{k \to \infty} \|T(z_{n_k})\| = \lim_{k \to \infty} 1 = 1
        \end{equation*}
        y donde hemos usado que toda sucesión acotada tiene una subsucesión convergente (Teorema de Bolzano-Weierstrass). \\

        \noindent
        Debido a la continuidad de $T^{-1}$ demostrada anteriormente:
        \begin{equation*}
            T^{-1}(T(z_{n_k})) \to T^{-1}(y) \implies z_{n_k} \to T^{-1}(y)
        \end{equation*}
        Sin embargo, como $z_n \to 0$, su subsucesión también debe converger a $0$. Esto implica que $T^{-1}(y) = 0$, y por la inyectividad de $T^{-1}$, $y = 0$. Esto contradice el hecho de que $\|y\| = 1$, concluyendo que $T$ es continuo.
    \end{itemize}
\end{proof}
\end{prop}

\begin{coro}
    Sea $(X, \|\cdot\|)$ un espacio normado de dimensión finita, entonces todas las normas sobre $X$ son equivalentes.
    \begin{proof}
    Sean $(X, \|\cdot\|_1)$ y $(X, \|\cdot\|_2)$ dos espacios normados con la misma estructura vectorial pero con normas diferentes, entonces
    \begin{equation*}
        (X, \|\cdot\|_1) \xrightarrow{T} \mathbb{R}^N \xrightarrow{T^{-1}} (X, \|\cdot\|_2)
    \end{equation*}
    con $T$ y $T^{-1}$ isomorfismos lineales topológicos. Por lo tanto, tenemos
    \begin{equation*}
        Id: (X, \|\cdot\|_1) \to (X, \|\cdot\|_2), \quad Id = T^{-1} \circ T
    \end{equation*}
    donde $Id$ es continua por composición de continuas, y entonces $Id$ está acotada, es decir, \mbox{$\|x\|_2\leq M\|x\|_1$}, $\forall x\in X$. Análogamente 
    \begin{equation*}
        Id: (X, \|\cdot\|_2) \to (X, \|\cdot\|_1), \quad Id = T \circ T^{-1}
    \end{equation*}
    es continua, y entonces $\|x\|_1 \leq C \|x\|_2 \ \forall x\in X$. Por lo tanto, las normas son equivalentes.
    
\end{proof}
\end{coro}



\begin{coro}
    Sea $(X, \|\cdot\|)$ un espacio normado de dimensión finita y $(Y, \|\cdot\|)$ un espacio normado cualquiera, entonces toda aplicación lineal
    \begin{equation*}
        L: X \longrightarrow Y
    \end{equation*}
    es continua.

\begin{proof}
    Como $X$ es de dimensión finita, entonces $L(X)=(\alpha_1, \ldots, \alpha_n)$ es un subespacio de dimensión finita en $Y$, de hecho podemos definir una base en dicho subespacio como $\{\beta_1, \dots, \beta_m\}$, con $m \le n$. Existe una matriz $M\in \mathbb{R}^{m \times n}$ tal que 
    $$L(x) = M[x]$$
    donde $[x]$ es el vector columna de las coordenadas de $x$ en la base de $X$. Tenemos entonces

    \begin{equation*}
        \|L(x)\| = \|M[x]\| \leq \|M\| \cdot \|[x]\| = C \cdot \|x\|
    \end{equation*}
    de donde obtenemos que $L$ está acotada, es decir, es continua.
\end{proof}
\end{coro}

\begin{coro}\label{cor:compacto-cerrado-acotado}
    Si $X$ es de dimensión finita, entonces 
    \begin{equation*}
        K\subseteq X \text{ es compacto} \iff K \text{ es cerrado y acotado}
    \end{equation*}
\begin{proof}
    \begin{description}
        \item[$\implies)$] Es cierto en cualquier espacio métrico.
        \item[$\impliedby)$] Sea $T: X \to \mathbb{R}^n$ un isomorfismo lineal topológico y $K \subseteq X$ cerrado y acotado. Entonces $T(K) \subseteq \mathbb{R}^n$ es compacto por ser $T$ continua, por lo que es cerrado y acotado en $\mathbb{R}^n$. Dado que $T^{-1}$ es continua, $K = T^{-1}(T(K))$ es compacto en $X$.
    \end{description}
\end{proof}
\end{coro}

\begin{ejercicio}
    Si $X$ es un espacio normado de dimensión finita, entonces $X$ es completo. \\ 
    \noindent
    Sea $T: X \to \mathbb{R}^n$ un isomorfismo lineal topológico. Si $(x_m) \subseteq X$ es una sucesión de Cauchy tenemos que $\|x_n - x_m\|_X \to 0$, entonces como $T$ es continua, y por tanto acotada, se tiene que

    $$\|T(x_n) - T(x_m)\|_{\mathbb{R}^n} \le \|T\| \cdot \|x_n - x_m\|_X$$
    como el lado derecho se va a cero, el lado izquierdo también, y entonces $(T(x_m)) \subseteq \mathbb{R}^n$ es una sucesión de Cauchy en $\mathbb{R}^n$, que es completo. Por lo tanto, $T(x_m) \to y \in \mathbb{R}^n$. Dado que $T^{-1}$ es continua, $x_m = T^{-1}(T(x_m)) \to T^{-1}(y) \in X$. Esto muestra que toda sucesión de Cauchy en $X$ converge, por lo que $X$ es completo.
\end{ejercicio}

\begin{ejercicio}
    Si $Y$ es un subespacio de dimensión finita en un espacio normado $X$, entonces $Y$ es cerrado. \\
    \noindent
    Sea $T: Y \to \mathbb{R}^n$ un isomorfismo lineal topológico. Dado que $\mathbb{R}^n$ es un espacio normado completo, es cerrado en sí mismo. Por lo tanto, $Y = T^{-1}(\mathbb{R}^n)$ es cerrado en $X$ por ser $T^{-1}$ continua.
\end{ejercicio}


\begin{lema}(de Riesz).
    Sea $X$ un espacio normado cualquiera y $Y \subsetneq X$ un subespacio propio, entonces

    \begin{equation*}
        \forall \varepsilon > 0 \quad \exists u_\varepsilon \in X : \|u_\varepsilon\| = 1 \quad \text{y} \quad d(u_\varepsilon, Y) \geq 1-\varepsilon
    \end{equation*}
    donde $d(u, Y) := \inf\limits_{y \in Y} d(u, y)$.
\begin{proof}
    Es importante destacar que estudiaremos solo el caso $0<\varepsilon<1$, porque para $\varepsilon >1$, la desigualdad $d(u_\varepsilon, Y) \geq 1-\varepsilon$ se cumple trivialmente para cualquier $u_\varepsilon$ con $\|u_\varepsilon\|=1$. \\

    \noindent
    Como $Y \subsetneq X$, entonces $\exists v \in X \setminus Y$ y por lo tanto $d(v, Y) = \delta > 0$. Sea $y_\varepsilon \in Y$ tal que
    $$\|v - y_\varepsilon\| \leq \frac{\delta}{1 - \varepsilon} > \delta$$
    \noindent
    Sea $u_\varepsilon = \dfrac{v - y_\varepsilon}{\|v - y_\varepsilon\|}$, entonces $\|u_\varepsilon\| = 1$ y es tal que para cualquier $z \in Y$:
    \begin{equation*}
        d(u_\varepsilon, z) = \|u_\varepsilon - z\| = \left\| \frac{v - y_\varepsilon}{\|v - y_\varepsilon\|} - z \right\| = \left\| \frac{v - y_\varepsilon - \|v - y_\varepsilon\| \cdot z}{\|v - y_\varepsilon\|} \right\|  \overset{(*)}{\geq} \frac{\delta}{\delta / (1 - \varepsilon)} = 1 - \varepsilon 
    \end{equation*}
    donde en $(*)$ hemos aplicado para el denominador la desigualda anterior y para el numerador, como sabemos que $y_{\varepsilon}$ y $z$ estan en $Y$ y es un subespacio, entonces es cerrado para la suma y producto de escalares, y tenemos

    \begin{equation*}
        \|v - y_\varepsilon - \|v - y_\varepsilon\|=\|v - (y_\varepsilon + \|v - y_\varepsilon\| \cdot z)\|=\|v-y^*\| \geq d(v, Y) = \delta \quad \text{ donde } y^*\in Y
    \end{equation*}
    lo que implica que $d(u_\varepsilon, Y) \geq 1 - \varepsilon$.
\end{proof}
\end{lema}

\begin{teo}\label{teo:compacto-dim-finita}
    Sea $X = (X, \|\cdot\|)$ un espacio normado, entonces:
    \begin{equation*}
        X \text{ es de dimensión finita} \iff \text{la bola unidad es compacta}
    \end{equation*}

\begin{proof}
    \begin{description}
        \item[$\implies)$] Ya está demostrado en Corolario \ref{cor:compacto-cerrado-acotado}.
        \item[$\impliedby)$] Sa $\bar{B} = \{x \in X \mid \|x\| = 1\}$ compacta. Demostremos que $X$ es de dimensión finita. \\ \\
        \noindent
        Sea $x_1 \in \bar{B}$ y sea $Y_1 = \langle x_1 \rangle$ el subespacio generado por $x_1$. Si $Y_1 = X$ hemos terminado. De lo contrario, $Y_1$ es cerrado y por Riesz, tenemos que
        \begin{equation*}
            \exists x_2 \in \bar{B} \text{ tal que } d(x_2, Y_1) > \frac{1}{2}
        \end{equation*}
        Definimos ahora $Y:= \langle x_1, x_2 \rangle$ el subespacio generado por $x_1$ y $x_2$. Si $Y = X$ hemos terminado, de lo contrario 
        \begin{equation*}
            \exists x_3 \in \bar{B} : d(x_3, Y) > \frac{1}{2}
        \end{equation*}
        y definimos $Y_3 := \langle x_1, x_2, x_3 \rangle$ e iteramos. Supongamos por reducción al absurdo que el proceso no termina, entonces
        \begin{equation*}
            \exists (x_n) \in \bar{B} \text{ tal que } d(x_n, x_i) > \frac{1}{2} \text{ para todo } n \in \{0, \dots, n-1\}
        \end{equation*}
        donde sabiendo que $X$ es espacio normado de dimensión finita, es completo y por tanto al no ser la sucesión $(x_n)$ de cauchy, no es convergente. De hecho, lo anterior afirma que $(x_n)$ no admite subsucesión convergente en $\bar{B}$, lo que contradice la compacidad de $\bar{B}$, y por reducción al absurdo se tiene.
    \end{description}
\end{proof}
\end{teo}

\begin{ejemplo}
    Demostremos que $l_2$ es de dimensión infinita. \\ 

    \noindent
    Sea $e_n=(0, 0, \dots, \overset{n}{1}, \dots, 0)$ la sucesión que tiene un $1$ en la posición $n$ y $0$ en las demás posiciones. Es fácil verificar que cada $e_n$ pertenece a $\bar{B} = \{x \in l_2 : \|x\|_2 = 1\}$, ya que $\|e_n\|_2 = 1$. Además, para cualquier $n \neq m$, se tiene:
    \begin{equation*}
        \|e_n - e_m\|_2 = \sqrt{(1-0)^2 + (0-1)^2} = \sqrt{2}
    \end{equation*}
    Esto implica que la sucesión $(e_n)$ no es de Cauchy, ya que la distancia entre cualquier par de términos es constante e igual a $\sqrt{2}$. Por lo tanto, $(e_n)$ no converge en $l_2$, entonces $\bar{B}$ no es compacto, y por el Teorema \ref{teo:compacto-dim-finita} tenemos que $l_2$ es de dimensión infinita.
\end{ejemplo}